\section{Previous Framework Versus the Current Framework}\label{sec:before-after}
The project started from a two-model design (a 3-head ResMLP and an ST-GCN) trained on a rule-derived table, with the rule engine deciding every disagreement. This section sets the previous framework against the current one, component by component. Every number is measured on photos or videos the model never trained on, with the same test set on both sides; the sources are Sections~\ref{sec:results} and~\ref{sec:compare} and \texttt{docs/BENCHMARKS.md}. Where the current framework is not better, that is stated in the last paragraph.

\begin{table}[t]
\centering\footnotesize
\caption{Previous versus current framework (held-out data only; bold marks the better value of each row).}\label{tab:before-after}
\begin{tabularx}{\textwidth}{@{}p{4.1cm}p{3.3cm}p{3.3cm}X@{}}
\toprule
Measure & Previous & Current & Test set \\ \midrule
Training labels & 654{,}488 rule-derived rows; 3-D angles in training, 2-D angles served (median mismatch 14\textdegree) & \textbf{cue-verified}: instructor names the pose, rules agree, pose is held; served 2-D angles & 24 videos, 1{,}206{,}391 frames \\
Overall naming accuracy, decision path in production & 36.9\% & \textbf{78.9\%} & 1{,}685 held-out public photos \\
False alarm (a pose outside the vocabulary named as one of ours), lower is better & 78.3\% & \textbf{20.8\%} & same 1{,}037 other-pose photos \\
Poses passing the 0.70 recall and precision bar & 0 & \textbf{7} & same photos \\
Naming on wild photos & 37.9\% & \textbf{55.3\%} & frozen Commons-103, out of fold \\
Score drops when one joint is broken by 45\textdegree & 59.2\% & \textbf{85.0\%} & 648 held-out photos $\times$ 5 trials \\
Score on correct form (higher is better) & 0.79 & \textbf{0.81} & same trials \\
ST-GCN macro recall on a rule-defined movement task (19 classes) & 14.2\% & \textbf{53.1\%} & 9{,}458 windows of the same 12 videos; the order-blind MLP scores 66.7\% \\
ST-GCN macro recall on held poses (the task the app's sequence row shows) & 18.8\% & \textbf{82.1\%} & 2{,}682 windows of 12 videos the old model never saw; outputs mapped to the 8-pose vocabulary \\
Share of true-pose windows the ST-GCN answers confidently (precision of those answers) & 2.8\% (93.8\%, 32 answers) & \textbf{50.5\%} (88.1\%, 611 answers) & same windows \\
ST-GCN macro recall at the training rate (25 frames/s) & 0.19 & \textbf{0.82} & same windows \\
Hidden joints & mirrored from the other side (median angle error 8.6\textdegree) & \textbf{visibility-aware: abstains rather than guesses}; 5.8\textdegree\ for the offline CLIFF experiment, not shipped & occlusion test, Section~\ref{sec:occ} \\
Poses with joint-level cues & 17 of 19, wrong joints marked red on Tree and Lunge & \textbf{17 of 19, no red on the two without bands} & live check on the deployed app \\
\bottomrule
\end{tabularx}
\end{table}

\runin{What produced the gain.} Three changes account for almost all of it. (i)~The labels: the previous table named a pose wherever the rules said so, so a model could at best copy the rules, and the rules won every disagreement; the new labels need the instructor's spoken pose name and the rules and a steady hold, and the features are computed the way the app serves them. (ii)~The decision path: the previous path let the rule engine overrule the model, so swapping in a better model changed nothing (identical 36.9\% and 78.3\%); the cascade lets a photo-trained model name the pose and a second model veto it when the pose is not one of ours, which is what removed most false alarms. (iii)~The ST-GCN data and input rate: the previous live model scored 18.8\% macro recall on held-out videos and was confident on under 3\% of windows; the browser also fed it frames at up to 2 per second instead of 25, which is a different input from what it was trained on. The current model is trained on cue-verified windows with mirror and photo-hold augmentation, and the web app now resamples every camera frame to 25 per second. At the old feed rate of about 2 frames per second neither model is usable (macro recall 0.48 for the previous model and 0.39 for the current one), so the buffer fix is what makes the temporal model usable at all.

\runin{Where the current framework is not better.} (a)~The per-joint deviation head is still weak: the broken joint is its first choice in 12.4\% of trials against 6.7\% by chance, so joint-level cues come from each pose's angle bands, not from that head. (b)~On wild Commons photos the previous MLP used alone names 60.2\% against 55.3\% for the cascade; the cascade gives up about five photos of naming accuracy for roughly 60 points fewer false alarms. (c)~The previous ST-GCN's few confident answers were more precise (93.8\% against 88.1\%), but they covered 2.8\% of the windows against 50.5\%. (d)~The previous ST-GCN was trained to name transitions: its encoder holds 30 classes (21 holds, seven directional transitions such as \texttt{transition:lunge\_pose->warrior\_2}, \texttt{transition:other} and \texttt{unrecognized}). The shipped model has one generic transition class, because the cue-verified labels cannot name transitions (Section~\ref{sec:flow}). The code comments and training notes quote 24 or 25 classes for this model and a macro recall of 63.0\% on two held-out videos with rule-defined labels; neither the class count nor that figure matches the checkpoint now deployed, so we do not use them.
(e)~A random forest trained on the same photos beats the whole cascade on same-source photos (85.3\% against 78.9\%) and loses on wild photos and on false alarms (Section~\ref{sec:results}); we judge the cascade by the wild and false-alarm columns because those are the conditions of a phone camera.

\subsection{The MLP, generation by generation}\label{sec:mlp-gens}
The pose-naming MLP has been retrained several times, and every generation is scored on the same held-out data: Table~\ref{tab:posehead} (naming and false alarms, from the two-month-old original to the current model), Table~\ref{tab:policy} (the decision path around it) and Table~\ref{tab:form} (does the form score react to a broken joint). In short: the original passed no pose and named a pose for 75.0\% of photos of poses we do not know; the first retrain was worse on wild photos (13.6\%) before the photo-domain data arrived; the 21 September model names wild photos best (60.2\%) but names a pose for 81.5\% of the unknown ones; the new gate model is the opposite (49.5\% on wild photos, 22.6\% false alarms, three poses passing); and the cascade built from the two passes seven poses with 20.8\% false alarms. No single MLP is best everywhere, which is why the current framework is a pair of small MLPs and not a replacement of one by another. The correctness score comes from the new gate model, which scores 0.81 on correct form and drops by 85.0\% when one joint is broken, against 0.79 and 59.2\% for the 21 September model.

\subsection{Is the sequence model a genuine ST-GCN?}\label{sec:genuine}
We checked this directly because a sequence model that is not a graph network but is called one would invalidate every sequence result. The evidence, from the repository and the training logs:
\begin{itemize}
\item \textbf{The code.} The model in \texttt{backend/app/models/sequence.py} applies, in each of three blocks, a $1{\times}1$ convolution followed by multiplication with the symmetric-normalised $33{\times}33$ adjacency of the MediaPipe skeleton (\texttt{einsum('vw,ncwt->ncvt')}), then a 9-tap temporal convolution with batch normalisation and a residual path.
\item \textbf{The parameter count.} Counting the layers by hand (blocks $3{\to}64{\to}128{\to}256$, head $256{\to}128{\to}9$) gives 893{,}897 parameters, which equals the count printed by the benchmark run for the 9-class model. A recurrent model of this size has a different count and different layer names.
\item \textbf{Strict loading.} The deployed loader and the old-versus-new benchmark load every checkpoint into this class with \texttt{load\_state\_dict}, which fails on any checkpoint with different layer names; all three previous checkpoints loaded and were scored.
\item \textbf{An ablation that changes the graph.} Replacing the adjacency with the identity matrix (no skeleton graph) changes macro recall from 84.8\% to 86.4\% over nine classes, so the graph is used by the network, though the gain is small (one seed).
\item \textbf{History.} The repository's first commit (10 June 2026) held a different model under the same class name: a $1$-D convolution, two residual GRU blocks and attention, with no skeleton graph. It was replaced by the graph network on 17 July 2026 (commit \texttt{018a7ba}), and the class name \texttt{YogaSequenceLSTM} was kept so existing checkpoints would load. Results for any checkpoint produced before the replacement would be results for that GRU model; this is why the old-versus-new comparison loads the old checkpoints with strict loading into the graph class.
\end{itemize}
\runin{Why this architecture.} The design notes record the split of work: the MLP checks one frame (pose, correctness, per-joint deviation) and the ST-GCN reads 60 frames for transitions and flow, because a single-frame classifier can judge alignment but not movement \cite{yan2018stgcn}. A graph network suits that job for three reasons: it models time (temporal convolution), the structure of the skeleton (graph convolution) and 3D, since it is the only model fed the full $(x,y,z)$ landmark coordinates while the MLP receives angles computed with $z$ set to zero because the camera's depth is the least reliable coordinate (Section~\ref{sec:features}). The project proposal fixed this pair of models; the proposal itself is not in the repository, so these reasons are taken from the design notes and from the team. What we measured for each reason: \emph{time}: on held poses a window-statistics MLP that ignores frame order matches the ST-GCN, and movement cannot be tested with our labels (Section~\ref{sec:flow}); \emph{graph}: removing the skeleton graph changes nothing measurable; \emph{3D}: the depth coordinate is the largest effect we measured (Table~\ref{tab:z-report}). Setting $z$ to zero lowers every sequence model on both tasks: on held poses the MLP on window statistics falls from 94.2\% to 59.3\% macro recall, the temporal CNN from 90.5\% to 70.1\% and the ST-GCN from 88.7\% to 28.6\%, where it collapses onto the generic transition class; on the movement task the losses are 13.2, 4.1 and 11.6 points. An older check on two held-out videos with rule-derived labels agrees in direction (ST-GCN macro recall 54.6\% with raw $z$, 49.7\% with bone-length-corrected $z$, 44.8\% with $z$ zeroed). So the project did not forget 3D: the sequence input is where it lives, and the per-frame MLPs deliberately use $z{=}0$ angles because that is what the browser serves and because the training angles had been computed with $z$ (a $14^\circ$ median gap, Section~\ref{sec:features}). The caveat is the signal itself: MediaPipe's $z$ is monocular relative depth, three ways of improving it (the library's world landmarks, a pretrained 2-D-to-3-D lifter, bone-length correction) did not improve on the baseline they were compared with in earlier internal experiments (world landmarks 42.9\% against 45.7\% for the 2-D angles on photos; the lifter 10\%; bone-length correction 49.7\% against 54.6\% for raw $z$ on two held-out videos), and the live $z$ from a phone camera may differ from the training videos; 3D is both the most informative and the least reliable input. A window-statistics MLP given the same $(x,y,z)$ window does as well as the ST-GCN, which suggests that the useful part of the sequence input is the 3-D information and not the graph or the temporal convolution; we have not yet tried giving the per-frame naming model that 3-D window.

\begin{table}[t]
\centering\footnotesize\setlength{\tabcolsep}{3pt}
\caption{Effect of the depth coordinate: macro recall of the same models trained with the full $(x,y,z)$ skeleton and with $z$ set to zero (bold: best in column). Held poses: nine classes on the cueT2 video folds; Movement: the 19-class rule-defined task. At most 3{,}000 training windows per class and one seed, so the full-skeleton held-pose values differ slightly from the sequence-baseline table of Section~\ref{sec:results}.}\label{tab:z-report}
\begin{tabular}{@{}lrrrr@{}}
\toprule
 & \multicolumn{2}{c}{Held poses} & \multicolumn{2}{c}{Movement} \\ \cmidrule(lr){2-3}\cmidrule(lr){4-5}
Model & $x,y,z$ & $z{=}0$ & $x,y,z$ & $z{=}0$ \\ \midrule
MLP on window mean/std & \textbf{94.2\%} & 59.3\% & \textbf{59.5\%} & 46.3\% \\
Temporal CNN & 90.5\% & \textbf{70.1\%} & 56.2\% & \textbf{52.1\%} \\
ST-GCN & 88.7\% & 28.6\% & 47.5\% & 35.9\% \\ \bottomrule
\end{tabular}
\end{table}


\subsection{The sequence model on movement}\label{sec:flow}
The ST-GCN exists to model how the body moves between poses, so held-pose accuracy (Table~\ref{tab:stgcn}, Section~\ref{sec:results}) is not the whole test. We tried to build a movement benchmark from our labels and report what the data allows.

\runin{What the labels support.} In the 24 videos the cue-verified labels contain 589 holds in 21 videos and 568 consecutive hold pairs, 205 of them between different poses. Within 150 frames (6\,s) there are 40 transitions between different poses, in 14 distinct directional pairs, and no pair occurs in three or more videos; within 500 frames there are 70 transitions in 33 pairs, and still none does. A named-transition benchmark with held-out videos is therefore not possible on cue-verified labels. The previous model's 30-class vocabulary, with seven named transitions, was built from rule-defined labels instead.

\runin{A rule-defined flow task.} We rebuilt that label scheme on the same videos: a window is a transition when its net angle change exceeds 15\textdegree/s and is named by the dominant pose of its first and last third (the pose of each frame is the cue-verified label where one exists and the rule engine's pose elsewhere); otherwise it is a hold, or \texttt{unrecognized}. Windows with more than 15\% ambiguous frames are dropped (28{,}544 usable windows). Keeping classes with at least 100 windows from at least three videos leaves 19 classes: 16 holds, \texttt{transition:other}, \texttt{unrecognized} and one named transition, cobra pose to downward dog (144 windows, 15 videos). So even the rule-defined scheme names a single transition on this data. The labels are defined by rules, so the task measures whether a network recovers rule-defined flow from the raw skeleton without the rules; it is not agreement with a human annotator.

\begin{table}[t]
\centering\footnotesize
\caption{Movement task (19 flow classes, rule-defined labels), previous versus current ST-GCN and the reference MLP, on the 9{,}458 windows of the 12 videos that arrived after the previous model was trained (clean for it). The current-architecture models are trained per fold and scored only on videos they never saw. Bold marks the best value in a column. Cobra to downward dog is the only named transition with enough support.}\label{tab:flow-oldnew}
\begin{tabular}{@{}lrrrr@{}}
\toprule
Model & Overall & Macro recall (19 classes) & Hold classes (16) & Cobra $\to$ down.\ dog recall \\ \midrule
Previous ST-GCN (30-class checkpoint) & 23.1\% & 14.2\% & 14.4\% & 0.0\% \\
ST-GCN, current architecture, trained on this task & 56.7\% & 53.1\% & 52.1\% & 78.3\% \\
MLP on window mean/std (order-blind) & \textbf{65.6\%} & \textbf{66.7\%} & \textbf{65.8\%} & \textbf{87.0\%} \\ \bottomrule
\end{tabular}
\end{table}

\runin{Previous against current.} Trained properly on cue-verified data, the current ST-GCN architecture is far better than the previous checkpoint on videos the previous one never saw: 53.1\% macro recall against 14.2\%, and it recognises the one named transition at 78.3\% where the previous model scores 0.0\%. This agrees with the held-pose result (82.1\% against 18.8\% macro recall, Table~\ref{tab:stgcn}). Two cautions: the label scheme is the one the previous model was trained on, which favours the previous model rather than the current one, and the previous model's other six named transitions have no test windows here, so its transition ability is not tested beyond cobra to downward dog. On all 24 videos (the previous model may have seen 12 of them) its macro recall is 20.9\%.

\runin{ST-GCN against simpler models.} The order-blind MLP is better than the ST-GCN on this task as well (66.7\% against 53.1\% on the clean subset). Under one recipe on all 24 videos (every fold scored on videos it never trained on), macro recall over the 19 classes is 58.6\% for the MLP, 56.8\% for the BiLSTM with attention, 55.6\% for the temporal CNN, 53.1\% for the LSTM, 48.4\% for the ST-GCN without the skeleton graph and 47.7\% for the ST-GCN: the ST-GCN is last, the skeleton graph again makes no measurable difference (48.4\% against 47.7\%) and the order-blind MLP is first. The order on the one named transition is the same (recall 68.1\% for the MLP, 50.7\% for the ST-GCN, 66.0\% for the LSTM).

\runin{Arrow of time.} To test whether any model uses frame order at all, every moving window was shown forward and time-reversed and the model had to say which. Chance is 50\%, and a model that ignores order must score exactly 50\%. Every model is at chance: 50.0\% for the MLP (exactly, as it must be), 51.1\% for the LSTM, 50.1\% for the BiLSTM, 49.9\% for the temporal CNN, 50.0\% for the ST-GCN without the graph and 50.0\% for the ST-GCN, and the training loss stays at $\ln 2=0.693$ for all of them except the LSTM, whose loss only begins to fall. The test is therefore uninformative: under this recipe (12 epochs, about 8{,}000 windows per fold) no model learns the direction of time, so it neither supports nor refutes the claim that the ST-GCN models movement.

\runin{Caveats.} Rule-defined labels; 24 videos, 12 of them in the clean subset; overlapping windows (stride 12 of 60), so the effective sample is the number of holds and videos; one seed; and the ST-GCN is the slowest model to train (about 290\,s per fold against 5--15\,s). What this section establishes is that the current ST-GCN is much better than the previous one; it does not establish that a graph network is the best choice for movement, because we found no task, with our labels, on which it beats a window-statistics MLP.
